\section{Discussion and limitations}
\subsection{Review should be evaluated by its direction, not its presence}
The advocate--critic result challenges the assumption that another reviewing pass is automatically safer. When the initial grader under-marks, a lower-only critic has no mechanism for restoring lost credit. The observed negative bias is consistent with that asymmetry. Advocate-only review was selected for dev MAE and bias, while the highest test QWK belongs to no review in the ablation. Maintaining the dev decision avoids retrospectively optimising the design to the held-out results.

This does not establish that critics are useless. A teacher might prefer stronger protection against unsupported awards in another setting, or a different critic might make better decisions. Here, removing the critic also weakens confidence separation because disagreements no longer trigger referral. Mark quality and review workload therefore have to be reassessed together. Neither the full pipeline's 62\% coverage nor the final configuration's 92\% coverage is a universal automation target.

\subsection{Limits of the data and reference marks}
SPRAG is typed and cannot validate handwriting recognition. Its human--human agreement is much higher than system agreement even after selection of the final configuration. The reference mean is a useful consensus, but it is not a proof of conceptual correctness. The question-disjoint split reduces leakage between answers to the same prompt; it does not remove shared course vocabulary, rubric conventions or other within-corpus similarities.

The handwritten set has only nine marked students and one teacher. Its question-clustered intervals reflect resampling of the observed questions, not a population study of teachers or institutions. The teacher's credit for short or vague responses differs from the system's rubric-bound practice. With no independent second grader, the study cannot determine whether these differences represent excessive teacher generosity, overly strict rubric interpretation or both. Low OCR confidence indicates uncertainty but does not quantify transcription accuracy.

The classroom offset is post hoc. Leave-one-student-out fitting limits direct reuse of the target student's marks, but does not establish transfer to another teacher, examination or recognition setting. Likewise, the SPRAG dev offset changes score location without improving the underlying interpretation of evidence. Offsets should remain visible analytic adjustments, rather than being presented as discovered rubric knowledge. The selected advocate-only default has not been evaluated on the handwritten set in this record.

\subsection{Security, feedback and unmeasured outcomes}
The injection result demonstrates an actual failure of prompt-only separation between instructions and answer data. The deterministic mitigation is wording-specific and was checked against the same attack formulation; other paraphrases may evade it. Flagging sends suspicious work to a teacher but does not guarantee recovery of an uninfluenced mark. The clean SPRAG check is not a general false-positive estimate for diverse classroom writing.

Feedback verification checks consistency with the grade, criterion and quote. It does not show that the advice is pedagogically helpful. Students have not rated the feedback, and no controlled rewrite or learning study was run. The Hindi example similarly demonstrates a working interface rather than measured multilingual accuracy. Claims about feedback effectiveness, fairness across Indian languages or the ability to assess diagrams would exceed the available evidence.

OCR character and word error rates, paired scan-versus-transcript grading, broad robustness perturbations and demographic or language-group analyses were not measured in the reported study. They are not represented by placeholder result tables. Future work should collect consented paired transcripts, independent classroom grading and student feedback ratings, then evaluate meaningful revisions under a controlled comparison. Confidence calibration and the teacher's actual review time also require dedicated measurements. Until then, the appropriate framing is selective auto-marking with inspectable teacher review, given the substantial gap to human--human agreement.
